\SourceLeaf{205}\begin{center}
\begin{tabular}{r|p{.82\linewidth}}
I & \\
 & Umgestellt\\
II & \\
 & An 7 geschraubt\\
III & \\
 & An 7 halb zurückgeschraubt\\
IV & \\
 & An 3 geschraubt, wenn zugleich die Libelle berichtigt werden soll.\\
V &
\end{tabular}
\end{center}

Die Schlüsse in umgekehrter Reihenfolge lauten also wie folgt.

(1) $v^4=0$, weil (2) $v^3=-v^2$ und kleine gleiche Bewegungen an 7 den Wert von $v$ um denselben Betrag ändern.

(2) $v^3=-v^2$, weil $v^2=u^1$, $v^3=-u^3$ und (3) $u^1=u^3$.

(3) $u^1=u^3$, weil $S^1=S^3=u$ von $u$ abhängt, also $u^1=u^3$.

Da hier also $S^1$ berücksichtigt wird, muss man offenbar schon vor I die Ebene $T$ in eine Lage gebracht haben, in der die Enden der Blase sichtbar sind.

Soll nun auch die Libelle berichtigt werden, so ist zu zeigen, dass $i^5=0$.

Dies erkennt man an der Normalstellung von $S$, die aber offenbar nicht unmittelbar mit $i=0$, sondern nur mit $g=0$ zusammenhängt.

Die Synthese geht nun am besten von den ursprünglichen Werten aus.
\[
S = M+h+i+u,
\]
\[
\Sigma = M-h+i+u,
\]
\[
\begin{array}{rccc}
 & u & v & i \\
\cline{2-4}
{[\mathrm{I}]} & v^1 & v^1 & i^1 \\
{[\mathrm{II}]} & -v^1 & v^1 & i^1 \\
{[\mathrm{III}]} & v^1 & -v^1 & i^1 \\
{[\mathrm{IV}]} & 0 & 0 & i^1 \\
{[\mathrm{V}]} & 0 & 0 & 0 \\
\end{array}
\]

Da $v$ bei 5 verschiedenen Stellungen nur 1\,mal $=+u$ und 4\,mal $=-u$ ist, kann man $v$ als Zeichen ganz außer Acht lassen und [nur] von $u$ sprechen.

